% Einheit 1 von 2 – Die Volumenformel (bank/rotationsvolumen/e1.jsonl, zusammenbau v0.1)
%% TODO Zeitmarke Einheit 1: Marken-Zeile nicht lesbar
%% TODO Zweigzeile Teil 1: was hier gelernt wird (Satz in Schülersprache aus dem Einheitstitel) – Einheit 1
\einheitenkopf[e1]{Einheit 1 von 2 · Die Volumenformel}
\zweigzeile{Abitur LK}
\setcounter{aufgabe}{7}

%% TODO Titel: Ich-kann-Satz fehlt in der Bank (Platzhalter: Kettenname) – Nr. 8
%% TODO Anweisung: ein Satz über der ersten Teilaufgabe fehlt in der Bank – Nr. 8
\begin{aufgabe}{Die Volumenformel}
%% TODO \rechenplatz{n}: Schrittzahl des Grundfalls fehlt in der Bank (nur bei fester Schreibform, 2.3 b)
\begin{teile}
\teil Was rotiert um welche Achse? Eine Schale steht aufrecht; ihr Längsschnitt ist die Parabel $y = 0{,}5x^2$, der Körper entsteht durch Drehung um die senkrechte Achse. Kreuze an, rechne nicht. \\ \kreuz{um die x-Achse – der Ansatz $\pi \cdot$ Integral über $(f(x))^2\,dx$ passt} \\ \kreuz{um die y-Achse – der Ansatz passt nicht, es wird über y integriert}
\teil $f(x) = x + 2$ rotiert über [0; 3] um die x-Achse – Rotationsvolumen? V ≈ \leerfeld
\teil Eine gläserne Leuchte ist rotationssymmetrisch; sie entsteht durch Rotation des Graphen von $f(x) = -\tfrac{1}{32}x^2 + 3$ zwischen den Randstellen −8 und 8 um die x-Achse; 1 LE = 1 cm. Wie viel Liter fasst sie? \hfill (FHR 2020) \\ \leerfeld[l]
\teil Ein Pflanzkübel besteht aus einem Kegelstumpf, der durch Rotation von $g(x) = 0{,}5x + 2$ über [0; 2] entsteht, und einem Zylinder mit dem Radius 3 über [2; 4]; beide rotieren um die x-Achse, 1 LE = 1 dm. Wie viel Liter Erde fasst er? \leerfeld[l]
\teil Die Seitenwand eines Behälters entsteht durch Rotation des Graphen von $q(x) = \sqrt{4x + 36}$, $0 \le x \le 12$, um die x-Achse; 1 LE = 1 cm, der Boden liegt bei $x = 0$. Auf dem Boden liegt eine Kugel mit dem Durchmesser 8 cm, sie ist vollständig unter Wasser. Weise nach, dass mehr als 1\,000 cm³ Wasser im Behälter sind. \hfill (Abitur 2018 LK)
\teil Die Seitenwand eines Behälters entsteht durch Rotation von $q(x) = \sqrt{3x + 25}$ um die x-Achse; 1 LE = 1 cm, der Boden liegt bei $x = 0$. Eine Kugel mit dem Durchmesser 6 cm liegt auf dem Boden und ist ganz unter Wasser. Gießt man 250 cm³ Wasser nach, steigt der Wasserstand um 1 cm. Stelle eine Gleichung für die Füllhöhe t vor dem Nachgießen auf; löse sie nicht. \hfill (Abitur 2018 LK)
\end{teile}
\end{aufgabe}

%% TODO Titel: Ich-kann-Satz fehlt in der Bank (Platzhalter: Kettenname) – Nr. 9
%% TODO Anweisung: ein Satz über der ersten Teilaufgabe fehlt in der Bank – Nr. 9
\begin{aufgabe}{Die Volumenformel – Fehler finden, Begründen, Anwendung}
\begin{teile}
\teil Finde den Fehler. $f(x) = x + 4$ rotiert über [0; 2] um die x-Achse. \rechnung{(f(x))^2 &= x^2 + 16 \\ V &= \pi \cdot (\tfrac{8}{3} + 32) = \tfrac{104}{3}\pi}
\teil Warum wird in der Volumenformel der Funktionsterm quadriert? Begründe.
\teil Ein Weinfass entsteht durch Rotation von $f(x) = 4 - 0{,}01x^2$ über [−10; 10] um die x-Achse; 1 LE = 1 dm. Reicht es für den Inhalt von 900 Flaschen zu je 0,75 l? \janein
\end{teile}
\end{aufgabe}
